\scanpage{316}\mspage{40}
4.4. Soit, plus généralement
\[
\mathcal F\subset\mathrm{Fl}(\mathrm{Cat})
\]
un ens. de flèches dans $\mathrm{Cat}$, en pensant
plus particulièrement au cas où $\mathcal F$ est un
des sept \deleted{catégories} \inserted{ens.} du diagramme
(4.2.1.) (p.~30) \deleted{\unsure{distinctes de|distincts de}} \inserted{autres que} $\mathrm{Fib}_W$.
\deleted{Supposons que $\mathcal F$ contienne les iso et
soit stable par changement de base.} Nous
pouvons formuler \deleted{des} \inserted{trois} propriétés similaires
pour $\mathcal F$ à celles formulées dans $\mathrm{Loc}(7,abc)$
dans le cas de $\mathcal F_W$, notons ces
propriétés
\[
A(\mathcal F),\ B(\mathcal F),\ C(\mathcal F),
\]
% slanted
\margin{\uline{NB} Ces conditions ont un sens si $\mathcal F\subset\mathrm{Fl}(M)$, si $(M,W)$ est une catégorie de modèles quelconque}
donc
\[
\left\{\begin{aligned}
A(\mathcal F):&\quad\mathcal F\subset\mathrm{Fib}_W\\
B(\mathcal F):&\quad\mathrm{Fl}(\mathcal F_S)\cap W_{\unsure{/S|S}}\subset W_S^u\\
&\quad\text{(où $\mathcal F_S$ désigne la sous-catégorie pleine de}\\
&\quad\text{$\mathrm{Cat}_{/S}$ formée des $f:X\text{---}S$ avec $f\in\mathcal F$)}\\
C(\mathcal F):&\quad\mathcal F\cap W\subset W^u.
\end{aligned}\right.
\tag*{(4.4.1)}
\]
Il est clair qu'on a $B(\mathcal F)\Longrightarrow C(\mathcal F)$, et je
vais \deleted{\unsure{montrer}} % Broken underline below both parts of the cancelled phrase.
\uline{\deleted{donner une condition sur $\mathcal F$ qui
implique}} qu'on a aussi $A(\mathcal F)\Longrightarrow B(\mathcal F)$,
% archive sheet number: 315
